% Status: result_ready_with_limits. Bounded-draft-ready.
% Evidence: interpretation_h2_daily_response
% Artefakte: h2_event_window_summary.csv; events_timeline_seed.csv
\chapter{H2: Ereignisfenster}
\label{ch:h2}

H2 prüft, ob sich Polymarket-Wahrscheinlichkeiten um sieben vorab kuratierte
öffentliche Ereignisse in plausibler Richtung und innerhalb fester Tagesfenster
bewegen \citep{lit_eventstudy_001,lit_emh_001}. Die Ereignisse wurden vor der
Analyse mit Zeitstempel und erwarteter Richtung fixiert.

\paragraph{Sichtbare Tagesreaktion.}
Die stärkste Primärfenster-Bewegung (0--1 Tag) zeigt das Attentat auf Trump am
13.~Juli 2024 mit einer abnormalen Veränderung von $+7.2$ Prozentpunkten zugunsten
von Trump. Die gerichtet erwarteten Ereignisse bewegen sich konsistent: die
Verurteilung im Schweigegeld-Prozess ($-4.1$ pp) und Bidens Rückzug zugunsten
von Harris ($-1.7$ pp für Trump) gehen in die erwartete Richtung. Die übrigen,
als neutral kuratierten Ereignisse zeigen kleinere Bewegungen. Insgesamt
reagieren die Polymarket-Wahrscheinlichkeiten auf Tagesebene plausibel auf die
kuratierten Ereignisse.

\begin{table}[ht]
\centering
\caption{H2 --- Abnormale Veränderung im Primärfenster (0--1 Tag) (T3)}
\label{tab:t3}
\small
\begin{tabular}{@{}llr@{}}
\toprule
Ereignis (2024) & erwartete Richtung & abn. Veränderung (pp) \\
\midrule
Trump-Verurteilung (30.05.) & Trump runter & $-4.1$ \\
Biden--Trump-Debatte (28.06.) & neutral & $+2.5$ \\
Attentat auf Trump (13.07.) & Trump rauf & $+7.2$ \\
Vance als VP (15.07.) & neutral & $+3.5$ \\
Biden-Rückzug (21.07.) & Harris rauf & $-1.7$ \\
Walz als VP (06.08.) & neutral & $+1.3$ \\
Harris--Trump-Debatte (11.09.) & neutral & $-2.9$ \\
\bottomrule
\end{tabular}
\end{table}

\begin{figure}[ht]
\centering
\includegraphics[width=0.85\textwidth]{thesis_h2_event_window_car.png}
\caption{Tägliche Ereignisfenster-Bewegungen (F2). Quelle:
\texttt{data/results/thesis\_h2\_event\_window\_car.png}.}
\label{fig:f2}
\end{figure}

\paragraph{Limitation.}
Die Analyse beruht auf Tagespreisen und kann keine Intraday-Reaktionszeit
identifizieren, H2 stützt damit eine Aussage über tägliche Reaktionsmuster,
nicht über sofortige Marktreaktion. Da die Ereignisse vorab fixiert wurden,
besteht keine nachträgliche Auswahl nach sichtbaren Preisbewegungen.
